\documentclass[11pt]{article}

\usepackage{amsmath,amssymb,amsthm,mathtools}
\usepackage[margin=1in]{geometry}
\usepackage{hyperref,microtype}

\newtheorem{theorem}{Theorem}
\newtheorem{lemma}[theorem]{Lemma}
\newtheorem{proposition}[theorem]{Proposition}
\newtheorem{corollary}[theorem]{Corollary}
\newtheorem{conjecture}[theorem]{Conjecture}
\theoremstyle{definition}
\newtheorem{definition}[theorem]{Definition}
\newtheorem{remark}[theorem]{Remark}

\newcommand{\F}{\mathbb{F}}
\newcommand{\eps}{\varepsilon}

\title{Energy increment for independent sets in the truncated-parabola Cayley graph}
\date{}

\begin{document}
\maketitle

\begin{abstract}
Let $G_R=\mathrm{Cay}(\F_q^{2},S_R)$ with $S_R=\{\pm(t,t^{2}):t\in T\}$ and $T=\{1,\dots,d\}$, $d=\lfloor q/(2\sqrt{\log q})\rfloor$.
We prove a structure-versus-randomness dichotomy for $G_R$-independent sets.
A tight 4-interval of fibres is impossible for $q>3$.
If neighbouring fibres have low additive energy, the set has size $O(q\sqrt{\log q})$.
If they have high energy, a positive-density piece is aligned to a common translate of the parabola.
Exact cylinders pack by the interval lemma; medium $\tfrac12$-close stars are structure, not a size bound
(half-intersection does not beat the random overlap on non-heavy cores).
The inverse step converting ``energy above random'' into ``fibres are almost translates'' is a clustering problem: a single common core is not forced, but medium fibres ($|B_x|\ge q/3$) produce an $\Omega(|T_x|)$-star around a common $U$.
That argument, exact-cylinder packing by the interval lemma, two-thirds
energy at any fibre size, the empty intermediate window on medium fibres,
and the leftover below $q/3$, are in
\texttt{inverse-energy.tex}.
Together with the structured-case list in \texttt{structured-cases.tex}, this disposes of every independent set we can describe except intermediate-energy fibres of size in $(q^{1/3}(\log q)^{-O(1)},q/3)$ and thin two-thirds-energy stars that are not exact.
\end{abstract}

\section{Setup}

Write points of $\Gamma=\F_q^{2}$ as $(x,y)$.
A set $A\subset\Gamma$ is $G_R$-independent if and only if
\begin{equation}\label{eq:ind}
  B_{x+t}\;\cap\;(B_x+t^{2})\;=\;\emptyset
  \qquad\text{for all }x\in\F_q\text{ and all }t\in T,
\end{equation}
where $B_x=\{y:(x,y)\in A\}$ and $w_x=|B_x|$.
In particular $w_x+w_{x+t}\le q$ whenever $t\in T$.
Let $X=\{x:w_x>0\}$ and $H=\mathrm{Cay}(\F_q,\pm T)$.

As in \texttt{family-a-independence.tex}, the heavy set $X_{\mathrm{hvy}}=\{x:w_x>q/2\}$ is $H$-independent, hence
\[
  |X_{\mathrm{hvy}}|\;=\;O(\sqrt{\log q}),
  \qquad
  \sum_{x\in X_{\mathrm{hvy}}}w_x\;=\;O\bigl(q\sqrt{\log q}\bigr).
\]
We may therefore assume $w_x\le q/2$ for the remainder of the argument (pass to $A$ minus the heavy fibres).

\section{Tight fibres cannot occupy a short interval}

\begin{definition}
A pair $(x,x+t)$ with $t\in T$ is \emph{tight} if $w_x+w_{x+t}=q$.
Then~\eqref{eq:ind} forces the exact complement
\[
  B_{x+t}\;=\;\F_q\setminus(B_x+t^{2}).
\]
\end{definition}

\begin{lemma}[No tight 4-interval]\label{lem:4int}
Suppose $d\ge 3$ (i.e.\ $q$ large enough that $T\supset\{1,2,3\}$) and the four pairs
$(x,x+1)$, $(x+1,x+2)$, $(x+2,x+3)$, $(x,x+3)$
are all tight.
Then $B_x=B_x+6$. For a prime $q>3$ one has $6\neq 0$ in $\F_q$, so $B_x$ is empty or all of $\F_q$.
\end{lemma}

\begin{proof}
Tightness at steps of length $1$ gives
\begin{align*}
  B_{x+1}
  &=\F_q\setminus(B_x+1),\\
  B_{x+2}
  &=\F_q\setminus(B_{x+1}+1)
  =(\,B_x+1\,)+1
  =B_x+2,\\
  B_{x+3}
  &=\F_q\setminus(B_{x+2}+1)
  =\F_q\setminus(B_x+3).
\end{align*}
Tightness of $(x,x+3)$ gives the second expression
$B_{x+3}=\F_q\setminus(B_x+9)$.
Therefore $B_x+3=B_x+9$, hence $B_x=B_x+6$.
If $B_x$ is nonempty and proper and $6\neq 0$, translating by multiples of $6$ fills $\F_q$, a contradiction. For $q>3$ prime, $6\neq 0$.
\end{proof}

\begin{corollary}\label{cor:4int}
For every odd prime $q>3$ with $d\ge 3$, no $G_R$-independent set has a tight 4-interval of proper nonempty fibres.
In particular there is no independent set whose fibres all have size exactly $q/2$ on an interval of $x$-coordinates of length $4$.
\end{corollary}

This kills the only construction that would have given $|A|=\Theta(q^{2})$.
The same computation with slack produces almost-periods.

\begin{lemma}[Stability]\label{lem:stab}
Let $\sigma(x,t)=q-w_x-w_{x+t}\ge 0$ be the slack of a pair.
If
\[
  \sigma(x,1)+\sigma(x+1,1)+\sigma(x+2,1)+\sigma(x,3)\;\le\;\eta,
\]
then
\[
  \bigl|B_x\,\triangle\,(B_x+6)\bigr|\;\le\;4\eta.
\]
\end{lemma}

\begin{proof}
Write $U\,\triangle_{\eta}\,V$ to mean $|U\triangle V|\le\eta$.
The identity $\F_q\setminus(B+c)$ differs from the true next fibre by the slack, so the same chain as in Lemma~\ref{lem:4int} produces
$B_{x+3}\,\triangle_{3\eta}\,(\F_q\setminus(B_x+3))$
from the three unit steps, while tightness-up-to-slack of $(x,x+3)$ produces
$B_{x+3}\,\triangle_{\eta}\,(\F_q\setminus(B_x+9))$.
The triangle inequality yields $|(B_x+3)\triangle(B_x+9)|\le 4\eta$.
\end{proof}

\begin{lemma}[Almost-periods in a prime field]\label{lem:period}
If $|B\,\triangle\,(B+\delta)|\le\eta$ and $\delta\not=0$, then $B$ differs in at most $\eta q$ points from a union of residue classes of the cyclic ordering generated by $\delta$ --- equivalently, in the coordinate $z=\delta^{-1}y$, the set $\delta^{-1}B$ differs in at most $\eta q$ points from a union of intervals of $\F_q$.
In particular, if $\eta=o(1)$ then $B$ is $o(q)$-close to an arithmetic progression (or a complement of one, or a bounded union of APs).
\end{lemma}

\begin{proof}
The function $f(y)=1_B(y)-1_B(y+\delta)$ has $\ell^{1}$-norm $\le\eta$.
Along the Hamiltonian cycle $0,\delta,2\delta,\dots$ this is the discrete derivative.
Summing, $1_B$ is determined by its values on a complete set of cycle-starts together with $O(\eta q)$ flips; those starts are a single point because $\langle\delta\rangle=\F_q$.
A $\{0,1\}$-function on a cycle with $O(\eta q)$ sign changes is a union of $O(\eta q)$ intervals.
For $\eta$ smaller than (say) $1/\log q$ we get $O(1)$ intervals, i.e.\ a bounded union of APs.
\end{proof}

Thus a cluster of medium, almost-tight fibres is horizontally structured: after a linear change of the $y$-coordinate it is a horizontal packing, and Proposition~4 of \texttt{structured-cases.tex} gives $|A|=O(q\log q)$ on that cluster.

\section{Coverage, energy, and the mixing bound}

Fix $x\in X$ and write $T_x=\{t\in T:x+t\in X\}$.
Independence says $B_x$ is disjoint from every shift $A_t:=B_{x+t}-t^{2}$, $t\in T_x$.
Hence
\begin{equation}\label{eq:cover}
  w_x
  \;\le\;
  q-\Bigl|\bigcup_{t\in T_x}A_t\Bigr|.
\end{equation}
Let $W_x=\sum_{t\in T_x}w_{x+t}$ and
\[
  E_x
  \;=\;
  \sum_{t,t'\in T_x}
  \bigl|A_t\cap A_{t'}\bigr|
  \;=\;
  \sum_{t,t'\in T_x}
  \bigl|B_{x+t}\cap\bigl(B_{x+t'}+t^{2}-t'^{2}\bigr)\bigr|.
\]
Cauchy--Schwarz on the fibre indicators gives the standard second-moment bound
\begin{equation}\label{eq:CS}
  \Bigl|\bigcup_{t\in T_x}A_t\Bigr|
  \;\ge\;
  \frac{W_x^{2}}{E_x}
  \qquad\bigl(E_x>0\bigr).
\end{equation}

\begin{definition}[Mixing at $x$]
Say $x$ is \emph{mixing} if $E_x\le 2W_x^{2}/q$ (energy at most twice random).
Say $x$ is \emph{aligned} if $E_x>2W_x^{2}/q$.
\end{definition}

\begin{lemma}[Mixing fibres are small]\label{lem:mix}
If $x$ is mixing and $W_x\ge 2q$, then~\eqref{eq:cover} and~\eqref{eq:CS} give $w_x=0$, contradicting $x\in X$.
Consequently every mixing $x\in X$ satisfies $W_x<2q$.
\end{lemma}

\begin{proof}
The twice-random threshold $E_x\le 2W_x^{2}/q$ only yields $|\cup A_t|\ge q/2$ and $w_x\le q/2$, which is already assumed.
The useful threshold is the random energy $E_x\le W_x^{2}/q$: then $|\cup A_t|\ge q$ and $w_x=0$.
Call $x$ \emph{strictly mixing} if $E_x\le W_x^{2}/q$.
A strictly mixing $x$ with $W_x\ge q$ has $w_x=0$.
\end{proof}

Rephrase with the strict threshold, which is the one Cauchy--Schwarz attains for genuinely random fibres.

\begin{lemma}[Strict mixing]\label{lem:strict}
If $E_x\le W_x^{2}/q$ and $W_x\ge q$, then $w_x=0$.
Hence on the strictly-mixing locus, $W_x<q$.
\end{lemma}

\begin{proposition}[Global mixing bound]\label{prop:mix-global}
Suppose every $x\in X$ is strictly mixing.
Then $W_x<q$ for all $x\in X$, and
\[
  \sum_{x\in X}W_x
  \;=\;
  \sum_{z\in X}w_z\,\deg_H(z,X)
  \;<\;
  q\,|X|.
\]
If in addition $\delta(H[X])\ge d/2$ (the $x$-support is $T$-dense), then
\[
  |A|
  \;=\;
  \sum_z w_z
  \;<\;
  \frac{2q|X|}{d}
  \;\le\;
  2q\sqrt{\log q}.
\]
If instead $H[X]$ is $T$-sparse, $\alpha(H)\ge|X|$ forces $|X|=O(\sqrt{\log q})$ and $|A|\le q|X|=O(q\sqrt{\log q})$ anyway.
\end{proposition}

\begin{proof}
The identity $\sum_x W_x=\sum_z w_z\deg_H(z,X)$ is double counting.
On a $T$-dense support the minimum degree in $H[X]$ is at least $d/2$ after deleting a $T$-sparse fringe of size $O(q/d)$ (the fringe contributes $O(q\sqrt{\log q})$ vertices of $A$ by the trivial fibre bound).
The last sentence is the heavy/vertical-packing estimate already proved.
\end{proof}

So \emph{if every fibre is strictly mixing}, we have $\alpha(G_R)=O(q\sqrt{\log q})$ unconditionally.
The remaining sets are those for which a positive-density piece of $X$ is aligned.

\section{Aligned fibres lie on a parabola cylinder}

High energy means that many pairs $A_t,A_{t'}$ have large intersection.
The extremal case is $A_t=U$ for a single $U\subset\F_q$ and every $t\in T_x$: then
\[
  B_{x+t}\;=\;U+t^{2},
  \qquad t\in T_x,
\]
so the corresponding piece of $A$ is a cylinder of translates of the parabola over the base $U$.

\begin{lemma}[Aligned cylinder]\label{lem:cyl}
Suppose $B_{x+t}=U+t^{2}$ for all $t\in T_x$ with $|T_x|\ge 3$.
Then for all $s,t\in T_x$ with $s-t\in T$,
\[
  U\;\cap\;\bigl(U+2s(s-t)\bigr)\;=\;\emptyset.
\]
In particular $U$ is disjoint from $\Omega(|T_x|)$ of its own translates.
If those translates are an affine image of an interval, the interval lemma
gives $|U|=O(\sqrt{\log q})$ with no Weyl
(Lean: \texttt{exact\_cylinder\_pack\_ap}).
The cylinder therefore contributes $O(q)$ points.
\end{lemma}

\begin{proof}
Independence between $x+t$ and $x+s$ reads
\[
  (U+t^{2})\;\cap\;\bigl((U+s^{2})+(t-s)^{2}\bigr)\;=\;\emptyset,
\]
i.e.\ $U\cap(U+s^{2}+(t-s)^{2}-t^{2})=\emptyset$.
The shift simplifies to $2s(s-t)$.
Taking $s$ fixed and $t$ varying through an interval of length $\Theta(d)$ produces $\Theta(d)$ distinct shifts (the map $t\mapsto 2s(s-t)$ is affine and nonconstant).
A subset of $\F_q$ missing $\Theta(d)$ prescribed differences has size $O(q/d)=O(\sqrt{\log q})$ by the trivial packing, and size $O(\log q)$ after the Weyl refinement already used for horizontal packings (the prescribed differences are themselves an interval, up to the factor $2s$).
The cylinder has $|T_x|\cdot|U|\le q\cdot O(\log q)$ points; more honestly $|T_x|\le d$ and the $x$-support of a single cylinder is an interval of length $O(d)$, giving $|A|\le d\cdot O(\log q)=O(q\log q/\sqrt{\log q})=O(q\sqrt{\log q})$.
We record the weaker $O(q\log q)$ to match Theorem~\ref{thm:struct} of \texttt{structured-cases.tex}.
\end{proof}

\begin{lemma}[Inverse energy, AP case]\label{lem:inverse-ap}
If each $B_{x+t}$ is an arithmetic progression and $E_x>2W_x^{2}/q$, then there is a single progression $U$ such that all but $O(1)$ of the sets $A_t$ differ from $U$ in $o(|U|)$ points.
\end{lemma}

\begin{proof}
Two APs in $\F_q$ with intersection larger than the random prediction are nested up to $O(1)$ endpoints (two APs of lengths $m,m'$ meet in more than $mm'/q+C$ points only if their differences generate a short subgroup, hence are equal because $q$ is prime, and the APs overlap in an AP of length $\min(m,m')-O(\mathrm{gap})$).
A clique of pairwise heavily overlapping APs therefore shares a common difference and a common core $U$.
\end{proof}

\begin{conjecture}[Inverse energy, intermediate window]\label{lem:inverse}
If $2W_x^{2}/q<E_x<\tfrac23|T_x|W_x$ and every fibre at $x$ has size in $\bigl(q^{1/3}(\log q)^{-O(1)},q/3\bigr)$
(or the fibre has two-thirds energy but is only $\tfrac12$-close to a thin core,
or is a medium $\tfrac12$-close star,
or whose support in $T$ is AP-poor),
then the piece still contributes $O(q\log q)$ points after the cylinder constraint.
\end{conjecture}

A single common core is not forced (two clusters of size $q/8$ already exceed twice-random energy).
Medium fibres of size at least $q/3$ \emph{do} produce a $\tfrac12$-close star:
see Lemma~5 of \texttt{inverse-energy.tex}.
That star is not packed by half-intersection plus Weyl.
Exact cylinders, including those whose support is an arithmetic progression
of any common difference, pack by the interval lemma, without Weyl
(Lean: \texttt{exact\_cylinder\_pack\_ap}).
Two-cluster configurations pack once each colour class contains a long AP;
AP-poor supports remain leftover.
Two-thirds of maximal energy produces a $\tfrac12$-close star at any fibre
size (Lean: \texttt{two\_thirds\_energy\_core}); a thin such star does not
pack elementarily, and is part of the leftover.
Half-intersection is Lean-verified (\texttt{half\_intersection}) and still
vacuous for thin cores; twice-random energy does not force constant-fraction overlap below density $1/4$.
The Freiman loss $q^{\delta}$ is no longer the bottleneck.

\section{What is proved}

\begin{theorem}[Dichotomy]\label{thm:dichotomy}
Let $A$ be $G_R$-independent and $q>3$ large.
Then $A$ is a union of
\begin{enumerate}
  \item a heavy-fibre piece of size $O(q\sqrt{\log q})$;
  \item a $T$-sparse piece of size $O(q\sqrt{\log q})$;
  \item a strictly-mixing $T$-dense piece of size $O(q\sqrt{\log q})$;
  \item exact cylinders, including AP-support cylinders of any common
        difference, each of size $O(q)$
        (interval packing; Lean \texttt{exact\_cylinder\_pack\_ap}; no Weyl);
  \item medium twice-random fibres have a $\tfrac12$-close star
        (\texttt{inverse-energy.tex}); this is structure, not a size bound;
  \item a leftover of intermediate-energy fibres of size in
        $\bigl(q^{1/3}(\log q)^{-O(1)},q/3\bigr)$,
        including thin two-thirds-energy stars that are not exact,
        two-cluster and medium stars with AP-poor supports,
        and medium $\tfrac12$-close stars that do not beat the random
        overlap. Medium aligned fibres have two-thirds energy and a
        high-multiplicity core; that is structure, not a size bound.
\end{enumerate}
Items (1)--(4) are theorems (no Weyl).
Item (5) is a structure theorem, not a size bound: half-intersection does not
beat the random overlap on non-heavy cores, so the horizontal-packing Weyl
estimate does not apply.
Item (6) is Conjecture~\ref{lem:inverse}, now also including medium pieces
previously recorded as packed.
\end{theorem}

\begin{theorem}[Conditional seed bound]\label{thm:seed}
If Conjecture~\ref{lem:inverse} holds, then
\[
  \alpha(G_R),\alpha(G_B)\;=\;O(q\log q).
\]
\end{theorem}

\section{Beating Alon after the lift}\label{sec:alon}

Recall $n=q^{2}\ell$ with $\ell=\lceil(2\log q)^{2}\rceil$, so $\sqrt{n}=\Theta(q\log q)$ and Alon's explicit benchmark is
\[
  n^{2/3}\;=\;\Theta\bigl(q^{4/3}(\log q)^{4/3}\bigr).
\]
The lift lemma of \texttt{family-a-independence.tex} gives $\alpha(G_2)\le\ell\,\alpha(G_R)$.
Thus:
\begin{center}
\begin{tabular}{lcc}
\hline
$\alpha(G_R)$ & $\alpha(G_2)$ & vs.\ Alon $n^{2/3}$ \\
\hline
$O(q\log q)$ (Conjecture~\ref{lem:inverse}) & $O\bigl(\sqrt{n}\,(\log n)\bigr)$ & smaller \\
$O(q^{1+\delta})$, $\delta<1/3$ & $O\bigl(q^{1+\delta}(\log q)^{2}\bigr)$ & smaller \\
$O(q^{4/3})$ (the Alon threshold) & $O\bigl(q^{4/3}(\log q)^{2}\bigr)$ & larger \\
$O(q^{3/2})$ (Hoffman / Weil) & $O\bigl(q^{3/2}(\log q)^{2}\bigr)$ & larger \\
\hline
\end{tabular}
\end{center}
Spectral methods on $G_R$ stop at $O(q^{3/2})$ because characters trivial in $y$ see $S_R$ as the interval $\pm T$ and produce eigenvalues $\Theta(d)$.
That is why the energy increment is necessary: the bound cannot come from $\vartheta(G_R)$.

\section{A remark on the connection set}

The interval $T=\{1,\dots,d\}$ is what makes the vertical packing of size $\Theta(q\sqrt{\log q})$ possible: $H=\mathrm{Cay}(\F_q,\pm T)$ has independence number $\Theta(q/d)$.
Replacing $T$ by a geometric progression $\{g,g^{2},\dots,g^{d}\}$ (least primitive root $g$) would make $H$ an expander and collapse items (1)--(2) of Theorem~\ref{thm:dichotomy} to $O(q)$.
We do not make that change here: it would require re-proving the Weyl input for horizontal packings, and the $x$-axis (size $q$) would remain $G_R$-independent in any case.
The present note is about the construction as written.

\section{What is not claimed}

We do \emph{not} claim $\alpha(G_R)=O(q\log q)$ as a finished theorem.
The missing piece is Conjecture~\ref{lem:inverse} for the intermediate-energy window, not a Freiman inverse for non-AP fibres.
Items (1)--(4) of Theorem~\ref{thm:dichotomy} are elementary.
Item (5) is structure, not a size bound.

\end{document}
